\documentclass[10pt,a4paper]{article}
\usepackage[margin=15mm,top=12mm,bottom=14mm]{geometry}
\usepackage{amsmath,amssymb}
\usepackage{booktabs}
\usepackage{tabularx}
\usepackage{enumitem}
\usepackage{microtype}
\usepackage[T1]{fontenc}
\usepackage[utf8]{inputenc}
\usepackage{lmodern}
\usepackage{xcolor}
\usepackage{tcolorbox}

\setlength{\parindent}{0pt}
\setlength{\parskip}{0.15em}

\begin{document}

% Header Block
\begin{center}
  {\Large\textbf{Master Thesis Supervisor Meeting -- Briefing Agenda}}\\[0.25em]
  {\large\textbf{Mode-I Phase-Field Adaptive Remeshing Benchmark Reproduction}}\\[0.2em]
  {\normalsize\textbf{Date:} Thursday, 17 September 2026 \quad $\cdot$ \quad \textbf{Duration:} 45 Minutes}\\[0.12em]
  {\footnotesize\textbf{Candidate:} Pruthviraja Reddy Vandavagali (Matr. Nr. 68865) \quad $\cdot$ \quad \textbf{Supervisors:} Prof. B.~Kiefer, Dr.~S.~Roth (IMFD)}
\end{center}
\vspace{-0.4em}
\hrule height 0.8pt
\vspace{0.4em}

\begin{tcolorbox}[colback=gray!5!white,colframe=gray!60!black,title=\textbf{Governing Supervisor Directive},top=2pt,bottom=2pt,left=4pt,right=4pt]
\centering
\small\emph{``We need to have understood everything related to the first model before we increase complexity.''}
\end{tcolorbox}

\vspace{0.2em}

\section*{Meeting Agenda \& Discussion Flow (45 Minutes)}
\vspace{-0.3em}

\begin{enumerate}[label=\textbf{\arabic*.},leftmargin=5mm,itemsep=3pt,topsep=2pt]
  \item \textbf{Mode-I Benchmark Definition \& Fixed-Mesh Reference Anchor (00--08 min)}
  \begin{itemize}[itemsep=1pt,topsep=1pt,leftmargin=4mm]
    \small
    \item Single-edge-notched square plate ($1.0 \times 1.0\,\mathrm{mm}$) with zero-gap sharp seam ($a_0 = 0.5\,\mathrm{mm}$).
    \item Fixed-mesh reference anchor (Job \texttt{1398090}, $15{,}192$ elements): $F_{\max} = 0.757778\,\mathrm{kN}$ at $u = 0.005857\,\mathrm{mm}$; authoritative unconstrained OLS structural stiffness $K_0 = 137.945520\,\mathrm{kN/mm}$ ($R^2 = 0.99999960$).
    \item Multi-quantity convergence criteria and UEL internal energy balance limitation note.
  \end{itemize}

  \item \textbf{Pandey--Kumar Native Adaptive Workflow \& Error Indicator (08--15 min)}
  \begin{itemize}[itemsep=1pt,topsep=1pt,leftmargin=4mm]
    \small
    \item Automated two-job procedure: coarse linear elastic pre-analysis $\to$ \texttt{MISESERI} extraction $\to$ native \texttt{adaptiveRemesh} $\to$ dual UEL/UMAT mesh generation $\to$ nonlinear fracture solve.
    \item Epistemic clarity on \texttt{MISESERI}: Abaqus stress discretization error indicator associated with recovered von Mises stresses in linear elasticity; not a phase-field or damage indicator.
    \item Canonical Mode-I pre-analysis provenance: $2{,}906$ finite elements, $2{,}906$ \texttt{WHOLE\_ELEMENT} scalar records (formal retirement of obsolete Mode-II $3{,}930$-row dataset).
  \end{itemize}

  \item \textbf{Resolution of the 71,320-Element Stiffness Defect (Priority Question A) (15--25 min)}
  \begin{itemize}[itemsep=1pt,topsep=1pt,leftmargin=4mm]
    \small
    \item Root-cause isolation: Abaqus \texttt{pre} 16-node card limit on free-format \texttt{*NSET} records without \texttt{GENERATE}. Abaqus issued warnings and retained only 16 of the 150 \texttt{N\_BOTTOM} nodes; 134 entries were deleted.
    \item Kinematic mechanism: unconstrained nodes permitted bottom-edge vertical lift up to $\approx 48.34\%$ of applied stroke, artificially softening structural stiffness down to $K_0 \approx 122.38\,\mathrm{kN/mm}$.
    \item Verified correction: wrapping records to $\le 16$ entries/line constrained all 150 bottom nodes to $u_2 = 0$, recovering exact stiffness: $K_0 = 138.021013\,\mathrm{kN/mm}$ (frozen, $+0.05\%$) and $K_0 \approx 137.821\,\mathrm{kN/mm}$ ($137.820804\,\mathrm{kN/mm}$, full fracture Job \texttt{1404933}, $-0.09\%$).
  \end{itemize}

  \item \textbf{Discrepancy Audit: 71,320 vs $\approx 13{,}941$ Finite Elements (Priority Question B) (25--35 min)}
  \begin{itemize}[itemsep=1pt,topsep=1pt,leftmargin=4mm]
    \small
    \item Multi-release invariance: Abaqus 2019 GA through 2023.HF4 generate $100.000\%$ bitwise identical $71{,}320$-element meshes with identical SHA-256 hashes.
    \item Exhaustive 15-factor OFAT audit: ruled out release shifts, load scale, output frequency, mesher controls, and multi-pass remeshing; tested boundary conditions and element formulations.
    \item Regional distribution: crack process corridor ($|y| \le 0.05\,\mathrm{mm}$) contains $9{,}061$ elements ($12.7\%$), while far-field and transition account for $62{,}259$ elements ($87.3\%$) of our reconstruction.
    \item Literature gap: accessible evidence does not identify which unpublished implementation detail accounts for the reported $\approx 13{,}941$ mesh; author input is required for definitive closure.
  \end{itemize}

  \item \textbf{Supervisor Decisions Required \& Next Steps (35--45 min)}
  \begin{itemize}[itemsep=1pt,topsep=1pt,leftmargin=4mm]
    \small
    \item \textbf{Option A (Recommended):} Accept Gate~5 as externally under-specified, retain $71{,}320$ as the publication-literal reconstruction, document $\approx 13{,}941$ as an unresolvable literature gap, and proceed with Mode-I thesis synthesis/documentation.
    \item \textbf{Option B:} Authorize transmission of the prepared 6-question author inquiry and keep Gate~5 open.
    \item \textbf{Explicit Hold:} Under both options, Gate~7 (ABAQUSER), Mode-II, and state transfer remain on hold.
  \end{itemize}
\end{enumerate}

\vfill
\hrule height 0.4pt
\vspace{0.2em}
{\scriptsize\textbf{Meeting Package Documentation:} Master 12-page Report (\texttt{SUPERVISOR\_MEETING\_REPORT\_MODE1\_2026-09-17.pdf}) $\cdot$ Provenance Ledger $\cdot$ Figures 1 \& 2 $\cdot$ Full Checksum Ledger}

\end{document}
